%\subsection{Etapas del procedimiento de inclusión y exclusión aplicado(actualmente esta en la seccion de arriba)}
%El procedimiento de inclusión y exclusión se diseñó con el objetivo de garantizar la relevancia y aplicabilidad de las herramientas encontradas al aplicar la metodología de búsqueda. El proceso se estructuró en tres etapas, cada una con criterios específicos que permitieron refinar de forma progresiva el conjunto inicial de resultados hasta obtener una selección final de soluciones a evaluar.

%\paragraph{Etapa 1 - Recuperación inicial de resultados}
%En la primera etapa, se ejecutaron las cadenas de búsqueda definidas sobre distintas fuentes de información. Estas incluyeron repositorios de código abierto como GitHub, motores de búsqueda convencionales como Google, marketplaces de software, además de aplicar herramientas de inteligencia artificial (Perplexity y Gemini) para complementar la búsqueda. El objetivo de esta fase fue obtener un conjunto amplio de herramientas potenciales capaces de cubrir los aspectos de interés definidos.

%\paragraph{Etapa 2 - Aplicación de criterios de inclusión}
%La segunda etapa correspondió a la aplicación de los criterios de inclusión, previamente definidos en la sección 2.1.2.2. En esta fase se filtraron las herramientas recuperadas, manteniendo únicamente aquellas que permitan acceder a la solución y mencionen en su documentación o presentación los requerimientos funcionales definidos. %Se comprobó que ofrecieran cobertura de los aspectos de interés definidos dentro del alcance del estudio, tales como los aspectos de interés definidos para el análisis entre requisitos y las historias de usuario generadas. Esta etapa permitió conformar un conjunto depurado de herramientas que cumplían con las condiciones necesarias para su posterior análisis empírico.

%\paragraph{Etapa 3 - Aplicación de criterios de exclusión y verificación final}
%Por último, la tercera etapa consistió en la aplicación de los criterios de exclusión. Se descartaron las herramientas que no pudieron ser ejecutadas o verificadas, aquellas que carecían de mantenimiento activo, las que no ofrecían mecanismos de verificación acordes a los aspectos de interés definidos y los prototipos conceptuales sin implementación funcional. También se eliminaron versiones equivalentes de una misma herramienta, priorizando las alternativas más completas. De este modo, el conjunto final de soluciones quedó compuesto únicamente por herramientas accesibles, activas y alineadas con los objetivos del estudio, garantizando una evaluación objetiva.

%En resumen, este procedimiento permitió pasar de un conjunto inicial amplio a una muestra final representativa de herramientas basadas en inteligencia artificial para la generación y análisis de historias de usuario, que sean capaces de cubrir los aspectos de interés definidos. Las herramientas seleccionadas fueron posteriormente analizadas en profundidad, aplicando los criterios definidos.